problem:  
Let \( M_i^n \subset S^{n+k} \) be a sequence of compact, smooth, **embedded minimal submanifolds** with  
\[
\int_{M_i} |A_{M_i}|^n \le C, \qquad \mathrm{Vol}(M_i)\le V_0,
\]
where \(A_{M_i}\) denotes the second fundamental form in the round sphere \(S^{n+k}\).  

Suppose the varifold measures \(V_i := \mathcal{H}^n\!\llcorner M_i\) converge weakly (as Radon measures) to a stationary integral varifold \(V_\infty\).  

**Question (Critical curvature compactness and curvature quantization):**  
Does the uniform \(L^n\)-bound on \(|A|\) imply that, after passing to a subsequence and suitable reparametrizations, the sequence \(M_i\) converges smoothly away from a finite set of singular points, and that at each singular point the curvature concentrates according to a **quantized curvature measure**
\[
|A_{M_i}|^n \, d\mu_{M_i} \rightharpoonup |A_{M_\infty}|^n\, d\mu_{M_\infty} + \sum_{j} \Theta_j \, \delta_{p_j}
\]
for some finite set of points \(\{p_j\}\subset S^{n+k}\) and universal quanta \(\Theta_j>0\)?  

Equivalently: in the critical integral curvature regime \(p=n\), does any failure of compactness for minimal submanifolds correspond to curvature “bubbling,” in the sense that blow-ups around concentration points yield complete minimal cones or finite-total-curvature minimal submanifolds in \(\mathbb{R}^{n+k+1}\)?  

Can one classify all possible bubble profiles arising in such \(L^n\)-bounded families, analogously to the energy quantization phenomena in harmonic map and Yang–Mills theory?

Why is it a "good" problem:  
This problem engages the **critical** analytic threshold \(p=n\), where existing elliptic and compactness techniques reach their natural limit. For \(p>n\), Sobolev–Morrey embeddings guarantee strong control; for \(p<n\), total curvature can disperse and produce wild singularities; but \(p=n\) is the delicate frontier where curvature concentration is possible yet potentially quantized.  

Exploring whether minimal submanifolds obey a quantization principle for curvature concentration—analogous to energy quantization for harmonic maps—would profoundly extend our understanding of geometric regularity at critical scales. It would build a deep bridge between geometric measure theory (varifold compactness), PDE regularity theory (elliptic estimates at borderline integrability), and global differential geometry (structure of minimal cones and bubbles).  

Its statement is simple but its implications vast: determining how curvature can concentrate under critical \(L^n\) control could lead to a new “bubble-tree” compactification theory for minimal submanifolds in spheres—one of the most fundamental unresolved directions at the analytic heart of minimal geometry.